\chapter{The Labeled Transition System and Context-Decorated HML}
\label{sec:lts-hml}
% ===================================================================

\section{Minimal Contexts and the LTS}

To obtain a bisimulation that is a \emph{congruence} for a GSLT, we follow
the approach of Milner, Sewell, and Leifer~\cite{leifer2000deriving}.
The key idea is to label transitions not merely by actions but by
\emph{minimal reactive contexts} --- the smallest environments needed to
trigger a given reduction.

\begin{definition}[Context]
A \emph{context} $K[-]$ for a GSLT $\GSLT$ is a term of $\terms(\GSLT)$
with a distinguished \emph{hole} $[-]$.  The application $K[P]$ substitutes
$P$ for the hole.  A context is \emph{minimal} for a term $P$ and rule $r$ if:
\begin{enumerate}[label=(\roman*)]
  \item $K[P] \rewrite_r Q$ for some $Q$;
  \item there is no proper sub-context $K'$ of $K$ satisfying (i).
\end{enumerate}
\end{definition}

The Milner--Sewell--Leifer construction produces, for any GSLT $\GSLT$, a
labeled transition system $\lts(\GSLT)$ in which transitions are labeled by
minimal contexts:
\[
  P \xrightarrow{K} Q \quad\text{iff}\quad K[P] \rewrite Q.
\]

Three refinements of that statement are needed before it can be leaned on, and
the book leans on it heavily. They are recorded here rather than where they bite.

\begin{remark}[Minimality is a universal property, and the ambient object goes]
\label{lts:rem:universal}
The subterm formulation above is the readable one, and it is not quite the
working one. Leifer and Milner define minimality by a universal property: a
transition $P \xrightarrow{K} d[r]$ is asserted when $l \rewrite r$ is a rule and
$d[l] = K[P]$ is an \emph{idempotent relative pushout} in the relevant slice
category, so that $K$ is a \emph{least} enabling context in the categorical sense
rather than merely one with no proper enabling subterm~\cite{leifer2000deriving}.
The two agree in easy cases and the universal property is what supports the
congruence proof.

There is a second departure. Leifer and Milner's reactive systems fix a
distinguished object as the domain of the pushouts, which in practice restricts
attention to ground terms. \cite{BIHOL} observes that the proofs do not depend on
it and adopts an \emph{open} version. That is what a GSLT needs, since a lambda
theory's terms carry variables and its contexts have interfaces rather than a
single ground sort.

The existence of the pushouts is not free. It is a hypothesis about the theory,
verified for particular theories rather than for GSLTs at large, and it is
listed among this book's unproved assumptions in the front matter.
\end{remark}

\begin{remark}[Labels carry processes, and are matched up to bisimilarity]
\label{lts:rem:higherorder}
In a higher-order setting a minimal enabling context still carries
\emph{processes}: for \rhoc\ the labels have the shape
$-\Par\mathrm{in}(n,\lambda x.q)$ with $q$ an arbitrary process. Matching such
labels syntactically yields a relation known to be too
fine~\cite{sangiorgi}. The remedy, which this book adopts throughout, is to
compare labels position by position, relating constructor skeletons by identity
and process payloads by the bisimilarity \emph{being defined}. The functional
whose greatest fixed point this is remains monotone --- the relation occurs only
positively --- so Knaster--Tarski still applies and the definition is sound.
\end{remark}

\begin{remark}[The congruence is relative to a class of contexts]
\label{lts:rem:relative}
It is usual to summarise the construction as ``the induced bisimulation is a
congruence''. Unqualified, that is false for the theory this book cares about
most. Congruence holds with respect to a \emph{class} $\mathcal{A}$ of admissible
contexts, and observations must be restricted to the same class. For \rhoc,
\cite{BIHOL} proves congruence under $\mathrm{out}(n,-)$, $\mathrm{in}(n,-)$,
$-\Par-$ and $\ast(@(-))$, and must exclude contexts with a hole \emph{under} the
quote, such as $\mathrm{out}(@(-),z)$. The reason is exactly reflection:
communication turns on name equality rather than on bisimilarity, so $p \bisim q$
with $p \neq q$ gives $@p \neq @q$, and the two contexts behave differently.

This is not a blemish to be apologised for. Making the admissible class a
parameter is what allows one theorem to cover both the congruence result for a
calculus and the comparison of two calculi, where the admissible contexts are the
image of the source's --- which is precisely the shape
Definition~\ref{ob:def:morphism} takes in Chapter~\ref{ch:catgslt}. The
restriction was discovered three times as an obstacle before it could be
recognised as a component.
\end{remark}

\section{Context-Decorated Hennessy--Milner Logic}

\begin{definition}[Context-Decorated HML]
\label{def:cdHML}
The formulae of \emph{context-decorated Hennessy--Milner logic}
$\HML(\ctx)$ are generated by:
\[
  \phi, \psi \;::=\; \top \;\mid\; \bot \;\mid\; \phi \wedge \psi \;\mid\;
  \neg\phi \;\mid\; \langle K \rangle \phi
\]
where $K$ ranges over contexts of $\GSLT$.
The satisfaction relation is:
\[
  u \sat \langle K \rangle \phi \quad\text{iff}\quad
  K[u] \rewrite u' \;\text{ in one step, and }\; u' \sat \phi.
\]
\end{definition}

\begin{theorem}[Adequacy]
For any GSLT $\GSLT$ equipped with the Milner--Sewell--Leifer LTS:
\[
  u \bisim v \quad\Longleftrightarrow\quad
  \forall \phi \in \HML(\ctx),\; u \sat \phi \Leftrightarrow v \sat \phi.
\]
\end{theorem}

\begin{remark}[Which logic this is adequate for]
\label{lts:rem:modalonly}
$\HML(\ctx)$ as defined above has modalities and boolean structure and nothing
else. The generated logics of Chapter~\ref{ch:hypercube} have in addition a
\emph{structural} layer, one connective per term constructor, and the theorem
above says nothing about it --- nor could it, since a decomposition is not a step.
The structural layer is strictly more discriminating than $\bisim$ in general, and
the resulting apparent conflict with adequacy is the subject of
Chapter~\ref{ch:obsext}. Everywhere this book says ``adequate'' of a generated
logic with structural connectives, the relation intended is the one that chapter
supplies.
\end{remark}

\section{The Logical Metric}

Fix an enumeration $(\phi_n)_{n \in \NN}$ of the formulae of $\HML(\ctx)$,
ordered compatibly with the partial order on contexts (smaller contexts
appearing earlier).

\begin{definition}[Logical Metric]
\[
  d_{\HML}(u, v) \;=\; 2^{-n}, \quad\text{where}\quad
  n = \min\bigl\{k \;\mid\; u \sat \phi_k \not\Leftrightarrow v \sat \phi_k\bigr\}.
\]
If $u \bisim v$, set $d_{\HML}(u, v) = 0$.
\end{definition}

\begin{proposition}
$d_{\HML}$ is an \emph{ultrametric} on $\terms(\GSLT)/{\bisim}$.
Terms differing only in behaviors triggered by large (expensive) contexts
are closer under $d_{\HML}$ than terms differing in behaviors triggered
by small (cheap) contexts.  This is the \emph{locality principle}: local
distinguishability weighs more than global.
\end{proposition}

% ===================================================================
